Meziroční nárůst indexu mzdových nákladů (Eurostat \ac{LCI}) podle odvětví \ac{NACE} v~ČR dokumentuje, že odvětvová divergence mzdového vývoje je v~ČR výraznější než v~zemích se silnými sektorovými \ac{KS}: odvětví s~těsnými trhy práce (IT, logistika) vykazují nárůsty výrazně nad celkovým průměrem, zatímco zemědělství a~pohostinství stagnují.
Tato polarizace má přímý vztah k~míře kolektivního pokrytí: odvětví, kde jsou uzavírány odvětvové nebo skupinové \ac{KS} (energetika, bankovnictví), vykazují stabilnější a~méně cyklické mzdové nárůsty, zatímco fragmentovaná odvětví kopírují hospodářský cyklus bez tlumícího koordinačního efektu.
Plná čára (celková podnikatelská sféra) zakrývá tuto heterogenitu a~může vytvářet falešný dojem koordinovaného vývoje; pohled na jednotlivá odvětví ukazuje, že „průměrný nárůst“ v~ČR je artefaktem statistické agregace, nikoli výsledkem institucionální koordinace.
Navrhované odvětvové \ac{KS} s~produktivitními ukazateli jako závaznou referencí pro sjednávání by transformovaly tento reaktivní vzorec do proaktivního, ve kterém mzdové nárůsty odrážejí produktivitní vývoj odvětví v~reálném čase.

Meziroční nárůst indexu mzdových nákladů (Eurostat \ac{LCI}) podle odvětví \ac{NACE} v~ČR dokumentuje, že odvětvová divergence mzdového vývoje je v~ČR výraznější než v~zemích se silnými sektorovými \ac{KS}: odvětví s~těsnými trhy práce (IT, logistika) vykazují nárůsty výrazně nad celkovým průměrem, zatímco zemědělství a~pohostinství stagnují.
Tato polarizace má přímý vztah k~míře kolektivního pokrytí: odvětví, kde jsou uzavírány odvětvové nebo skupinové \ac{KS} (energetika, bankovnictví), vykazují stabilnější a~méně cyklické mzdové nárůsty, zatímco fragmentovaná odvětví kopírují hospodářský cyklus bez tlumícího koordinačního efektu.
Plná čára (celková podnikatelská sféra) zakrývá tuto heterogenitu a~může vytvářet falešný dojem koordinovaného vývoje; pohled na jednotlivá odvětví ukazuje, že „průměrný nárůst“ v~ČR je artefaktem statistické agregace, nikoli výsledkem institucionální koordinace.
Navrhované odvětvové \ac{KS} s~produktivitními ukazateli jako závaznou referencí pro sjednávání by transformovaly tento reaktivní vzorec do proaktivního, ve kterém mzdové nárůsty odrážejí produktivitní vývoj odvětví v~reálném čase.

Meziroční nárůst indexu mzdových nákladů (Eurostat \ac{LCI}) podle odvětví \ac{NACE} v~ČR dokumentuje, že odvětvová divergence mzdového vývoje je v~ČR výraznější než v~zemích se silnými sektorovými \ac{KS}: odvětví s~těsnými trhy práce (IT, logistika) vykazují nárůsty výrazně nad celkovým průměrem, zatímco zemědělství a~pohostinství stagnují.
Tato polarizace má přímý vztah k~míře kolektivního pokrytí: odvětví, kde jsou uzavírány odvětvové nebo skupinové \ac{KS} (energetika, bankovnictví), vykazují stabilnější a~méně cyklické mzdové nárůsty, zatímco fragmentovaná odvětví kopírují hospodářský cyklus bez tlumícího koordinačního efektu.
Plná čára (celková podnikatelská sféra) zakrývá tuto heterogenitu a~může vytvářet falešný dojem koordinovaného vývoje; pohled na jednotlivá odvětví ukazuje, že „průměrný nárůst“ v~ČR je artefaktem statistické agregace, nikoli výsledkem institucionální koordinace.
Navrhované odvětvové \ac{KS} s~produktivitními ukazateli jako závaznou referencí pro sjednávání by transformovaly tento reaktivní vzorec do proaktivního, ve kterém mzdové nárůsty odrážejí produktivitní vývoj odvětví v~reálném čase.

Meziroční nárůst indexu mzdových nákladů (Eurostat \ac{LCI}) podle odvětví \ac{NACE} v~ČR dokumentuje, že odvětvová divergence mzdového vývoje je v~ČR výraznější než v~zemích se silnými sektorovými \ac{KS}: odvětví s~těsnými trhy práce (IT, logistika) vykazují nárůsty výrazně nad celkovým průměrem, zatímco zemědělství a~pohostinství stagnují.
Tato polarizace má přímý vztah k~míře kolektivního pokrytí: odvětví, kde jsou uzavírány odvětvové nebo skupinové \ac{KS} (energetika, bankovnictví), vykazují stabilnější a~méně cyklické mzdové nárůsty, zatímco fragmentovaná odvětví kopírují hospodářský cyklus bez tlumícího koordinačního efektu.
Plná čára (celková podnikatelská sféra) zakrývá tuto heterogenitu a~může vytvářet falešný dojem koordinovaného vývoje; pohled na jednotlivá odvětví ukazuje, že „průměrný nárůst“ v~ČR je artefaktem statistické agregace, nikoli výsledkem institucionální koordinace.
Navrhované odvětvové \ac{KS} s~produktivitními ukazateli jako závaznou referencí pro sjednávání by transformovaly tento reaktivní vzorec do proaktivního, ve kterém mzdové nárůsty odrážejí produktivitní vývoj odvětví v~reálném čase.

\input{texparts/python/rscp_sector_lci_growth}



